\begin{monsheet}
\addcontentsline{toc}{section}{\ExamTocProblemLabel}

\begin{daimon}{1}
次の各問いに答えよ。\hfill（配点 20点）
\begin{shomon}
  \item $x=\dfrac{\sqrt{6}+\sqrt{2}}{\sqrt{6}-\sqrt{2}}$，$y=\dfrac{\sqrt{6}-\sqrt{2}}{\sqrt{6}+\sqrt{2}}$ とするとき，$x+y$，$xy$ の値をそれぞれ求めよ。さらに，$x^{3}+y^{3}$ および $x^{5}+y^{5}$ の値を求めよ。
  \item 円に内接する四角形 $\Pt{ABCD}$ において，$\Pt{AB}=3$，$\Pt{BC}=4$，$\Pt{CD}=5$，$\Pt{DA}=6$ とする。$\cos\angle\Pt{ABC}$ の値と四角形 $\Pt{ABCD}$ の面積 $S$ を求めよ。また，対角線 $\Pt{AC}$ の長さを求めよ。
  \item 硬貨 $\Pt{A}$ は表が出る確率が $\dfrac{2}{3}$，硬貨 $\Pt{B}$ は表が出る確率が $\dfrac{1}{4}$ である。まず，さいころを $1$ 回投げて，$1$ または $2$ の目が出たら硬貨 $\Pt{A}$ を選び，$3$ 以上の目が出たら硬貨 $\Pt{B}$ を選ぶ。次に，選んだ硬貨を $3$ 回続けて投げる。表がちょうど $2$ 回出る確率 $p$ を求めよ。また，表がちょうど $2$ 回出たとき，選ばれた硬貨が $\Pt{A}$ であった条件付き確率 $q$ を求めよ。
  \item 不等式 $\log_{3}(x-1)+\log_{3}(7-x)<2$ を満たす実数 $x$ の値の範囲を求めよ。また，$\log_{10}2=0.3010$，$\log_{10}3=0.4771$ として，$6^{40}$ は何桁の整数か求め，さらに $6^{40}$ の最高位の数字を求めよ。
\end{shomon}
\end{daimon}

\begin{daimon}{2}
実数 $x$ の関数 $f(x)=x^{3}-3x^{2}-9x+11$ を考える。また，$xy$ 平面上の曲線 $y=f(x)$ を $C$ とする。次の問いに答えよ。\hfill（配点 20点）
\begin{shomon}
  \item 関数 $y=f(x)$ の増減と極値を調べ，曲線 $C$ の概形をかけ。また，曲線 $C$ と $x$ 軸の共有点の座標をすべて求めよ。
  \item すべての実数 $h$ に対して $f(1+h)+f(1-h)=0$ が成り立つことを示し，曲線 $C$ は点 $\Pt{A}(1, 0)$ に関して対称であることを示せ。また，点 $\Pt{A}$ における曲線 $C$ の接線 $\ell$ の方程式を求めよ。
  \item 曲線 $C$ の極大点を $\Pt{P}$ とし，点 $\Pt{P}$ における曲線 $C$ の接線を $m$ とする。直線 $m$ と曲線 $C$ の共有点のうち点 $\Pt{P}$ 以外の点を $\Pt{Q}$ とするとき，点 $\Pt{Q}$ の座標を求め，曲線 $C$ と直線 $m$ で囲まれた部分の面積 $S_{1}$ を求めよ。
  \item 実数 $t$ が $-1<t<5$ の範囲を動くとき，直線 $x=t$ と直線 $m$ の交点を $\Pt{M}$，直線 $x=t$ と曲線 $C$ の交点を $\Pt{N}$ とする。線分 $\Pt{MN}$ の長さが最大となるときの $t$ の値と最大値を求めよ。また，そのときの直線 $x=t$ によって (3) の図形（面積 $S_{1}$）が $2$ つの部分に分けられるとき，$x\leqq t$ の部分の面積 $S_{2}$ と $S_{1}$ の比 $S_{2}:S_{1}$ を最も簡単な整数の比で表せ。
\end{shomon}
\end{daimon}

\begin{daimon}{3}
実数 $\theta$ に対して $t=\cos\theta$ とおく。次の問いに答えよ。\hfill（配点 20点）
\begin{shomon}
  \item 加法定理を用いて $\cos 3\theta$ を $t$ の多項式で表せ。また，方程式 $\cos 3\theta=\dfrac{1}{2}$ を満たす $\theta$（$0<\theta<\pi$）をすべて求めよ。
  \item $3$ 次方程式 $8t^{3}-6t-1=0$ の $3$ つの実数解を，(1) を用いて $\cos\theta$ の形で表せ。また，それらを小さい順に $\alpha, \beta, \gamma$（$\alpha<\beta<\gamma$）とするとき，$\alpha+\beta+\gamma$，$\alpha\beta+\beta\gamma+\gamma\alpha$，$\alpha\beta\gamma$ の値，および $\alpha^{2}+\beta^{2}+\gamma^{2}$ の値をそれぞれ求めよ。
  \item 関数 $g(t)=8t^{3}-6t-1$ の増減と極値を調べ，(2) の $3$ 解 $\alpha, \beta, \gamma$ がそれぞれ $-\dfrac{\sqrt{3}}{2}<\alpha<-\dfrac{1}{2}$，$-\dfrac{1}{2}<\beta<0$，$\dfrac{\sqrt{3}}{2}<\gamma<1$ を満たすことを示せ。
  \item 実数の定数 $a$ に対し，$0\leqq\theta\leqq\pi$ における $\theta$ の方程式 $\cos 3\theta-3a\cos\theta=0$ の異なる実数解 $\theta$ の個数を，$a$ の値によって分類せよ。
\end{shomon}
\end{daimon}

\clearpage

\begin{daimon}{4}
数列 $\{a_{n}\}$ の初項から第 $n$ 項までの和を $S_{n}$ とする。すべての自然数 $n$ に対して
\[
  S_{n} = \frac{3}{2}a_{n} - 2^{n+1} + 3
\]
が成り立つとする。次の問いに答えよ。\hfill（配点 20点）
\begin{shomon}
  \item $a_{1}$，$a_{2}$，$a_{3}$ の値をそれぞれ求めよ。また，$n\geqq 1$ のとき $a_{n+1}$ を $a_{n}$ と $n$ を用いて表せ。
  \item $b_{n}=\dfrac{a_{n}}{2^{n}}$ とおくとき，$b_{n+1}$ を $b_{n}$ を用いて表せ。また，数列 $\{b_{n}\}$ および数列 $\{a_{n}\}$ の一般項をそれぞれ $n$ の式で表せ。
  \item 初項から第 $n$ 項までの和 $S_{n}$ を $n$ の式で表せ。また，$a_{n}>10\cdot 2^{n}$ を満たす最小の自然数 $n$ を求めよ。
  \item 自然数 $n$ に対して $c_{n}=\dfrac{a_{n}+2^{n+2}}{10}$ とおくとき，和 $\dsp W_{n}=\sum_{k=1}^{n}\frac{k}{c_{k}}$ を $n$ の式で表せ。
\end{shomon}
\end{daimon}

\begin{daimon}{5}
四面体 $\Pt{OABC}$ において，$\Pt{OA}=2$，$\Pt{OB}=\Pt{OC}=\sqrt{5}$，$\Pt{AB}=\Pt{AC}=\sqrt{5}$，$\Pt{BC}=2$ であるとする。$\vec{a}=\Vec{OA}$，$\vec{b}=\Vec{OB}$，$\vec{c}=\Vec{OC}$ とおく。次の問いに答えよ。\hfill（配点 20点）
\begin{shomon}
  \item 内積 $\vec{a}\cdot\vec{b}$，$\vec{a}\cdot\vec{c}$，$\vec{b}\cdot\vec{c}$ の値をそれぞれ求めよ。また，$\triangle\Pt{OBC}$ の面積を求めよ。
  \item 頂点 $\Pt{A}$ から平面 $\Pt{OBC}$ に下ろした垂線の足を $\Pt{H}$ とする。$\Vec{OH}$ を $\vec{b}$，$\vec{c}$ を用いて表し，線分 $\Pt{AH}$ の長さと四面体 $\Pt{OABC}$ の体積 $V$ を求めよ。
  \item 辺 $\Pt{OA}$ の中点を $\Pt{M}$，辺 $\Pt{BC}$ の中点を $\Pt{N}$ とする。直線 $\Pt{MN}$ は辺 $\Pt{OA}$ と辺 $\Pt{BC}$ の両方に垂直であることを示し，線分 $\Pt{MN}$ の長さを求めよ。
  \item $4$ 点 $\Pt{O}$，$\Pt{A}$，$\Pt{B}$，$\Pt{C}$ をすべて通る球面（四面体 $\Pt{OABC}$ の外接球）の中心を $\Pt{K}$ とする。$\Vec{OK}$ を $\vec{a}$，$\vec{b}$，$\vec{c}$ を用いて表し，この外接球の半径 $R$ を求めよ。
\end{shomon}
\end{daimon}

\scratchpage
\end{monsheet}
